% part: intuitionistic-logic
% chapter: propositions-as-types
% section: proofs-to-terms

\documentclass[../../../include/open-logic-section]{subfiles}

\begin{document}

\olfileid[id]{int}{pty}{pt}

\olsection{Mengonversi \usetoken{P}{derivation} Menjadi Term Bukti}

Proses mengonversi !!{proof}s deduksi alami menjadi term bukti, yakni
menerjemahkan !!{proof}s dalam~\Log{N2i} menjadi !!{proof}s
dalam~\Log{tN2i}, sangat langsung. Kita cukup menulis suatu term bukti di
sebelah kiri setiap !!{formula} yang berada di sebelah kanan sekuen dalam
!!{derivation}, sehingga diperoleh sekuen berbentuk $\Gamma \Sequent M :
!A$. Kemudian kita mengatakan bahwa $M$ \emph{menyaksikan}~$!A$ dalam
konteks~$\Gamma$. Term bukti yang harus ditulis sepenuhnya ditentukan oleh
aturan-aturan \Log{tN2i}.

\begin{prop}
Setiap !!{proof}~$\delta$ atas $\Gamma \Sequent !A$ dalam \Log{N2i}
bersesuaian dengan !!a{proof}~$\delta'$ atas $\Gamma \Sequent N : !A$
dalam \Log{tN2i}. Term bukti $N$ ditentukan secara unik oleh~$\delta$.
\end{prop}

\begin{proof}
Kita melakukan induksi pada tinggi~$\delta$. Jika $\delta$ merupakan
aksioma $x:!A \Sequent !A$, maka $\delta'$ merupakan aksioma
$x:!A \Sequent x:!A$. Jika $\delta$ berakhir dengan suatu aturan inferensi,
hipotesis induksi memberikan term bukti bagi setiap premis serta
\Log{tN2i}-!!{proof}s atas sekuen yang memuat term bukti tersebut dan
bersesuaian dengan setiap premis. Kita menerapkan aturan \Log{tN2i} yang
bersesuaian; term bukti yang bersesuaian ditentukan oleh aturan inferensi
\Log{tN2i} itu.
\end{proof}

\begin{ex}
  Tinjau !!{proof} yang sangat sederhana dalam \Log{N1i} atas
  $(!A \land !B) \lif (!B \land !A)$:
  \[
  \AxiomC{$\Discharge{!A\land !B}{x}$}
  \RightLabel{\Elim\land[2]}
  \UnaryInfC{$!B$}
  \AxiomC{$\Discharge{!A\land !B}{x}$}
  \RightLabel{\Elim\land[1]}
  \UnaryInfC{$!A$}
  \RightLabel{\Intro\land}
  \BinaryInfC{$!B \land !A$}
  \DischargeRule{\Intro\lif}{x}
  \UnaryInfC{$(!A \land !B) \lif (!B \land !A)$}
  \DisplayProof
  \]
  Dalam \Log{N2i}, bukti tersebut tampak seperti ini:
  \[
  \Axiom$x : !A\land !B \fCenter !A\land !B$
  \RightLabel{\Elim\land[2]}
  \UnaryInf$x : !A\land !B \fCenter !B$
  \Axiom$x : !A\land !B \fCenter !A\land !B$
  \RightLabel{\Elim\land[1]}
  \UnaryInf$x : !A\land !B \fCenter !A$
  \RightLabel{\Intro\land}
  \BinaryInf$x : !A\land !B \fCenter !B \land !A$
  \RightLabel{\Intro\lif}
  \UnaryInf$\fCenter (!A \land !B) \lif (!B \land !A)$
  \DisplayProof
  \]
  Kita memberikan term bukti dari atas ke bawah. Term bukti bagi aksioma
  sama dengan variabel dalam konteks, yakni~$x$. Term bukti yang terkait
  dengan $\Elim\land[i]$ kemudian ditentukan sebagai $\proj{2}{x}$
  dan~$\proj{1}{x}$. Ketika menerapkan \Intro{\land}, kita menggabungkan
  kedua term bukti ini: $\pair{\proj{2}{x}}{\proj{1}{x}}$. Terakhir,
  aturan \Intro{\lif} menerapkan abstraksi-$\lambd$, mengikat~$x$ dan
  mencatat bahwa $x$ melabeli asumsi~$!A \land !B$:
  $\lambd[\typeof{x}{!A \land
  !B}][\pair{\proj{2}{x}}{\proj{1}{x}}]$.
  !!{proof} dalam \Log{tN2i} kemudian adalah:
  \[
  \Axiom$x : !A\land !B \fCenter x: !A\land !B$
  \RightLabel{\Elim\land[2]}
  \UnaryInf$x : !A\land !B \fCenter \proj{2}{x} : !B$
  \Axiom$x : !A\land !B \fCenter x: !A\land !B$
  \RightLabel{\Elim\land[1]}
  \UnaryInf$x : !A\land !B \fCenter \proj{1}{x} : !A$
  \RightLabel{\Intro\land}
  \BinaryInf$x : !A\land !B \fCenter \pair{\proj{2}{x}}{\proj{1}{x}} :
    !B \land !A$
  \RightLabel{\Intro\lif}
  \UnaryInf$\fCenter \lambd[\typeof{x}{!A \land
  !B}][\pair{\proj{2}{x}}{\proj{1}{x}}] : (!A \land !B) \lif (!B \land !A)$
  \DisplayProof
  \]
\end{ex}

\end{document}
